\section*{Abstract}
\addcontentsline{toc}{section}{Abstract}

This thesis studies a rolling-horizon warehouse-distribution planning problem motivated by last-mile delivery systems with flexible service days, gradually released orders, multi-trip vehicles, and depot-side execution limits. The core challenge is that high service quality cannot be achieved by static route planning alone: daily admission decisions affect future feasibility, while warehouse resources such as picking capacity, loading gates, and staging space can become locally overloaded even when route plans appear transportation-feasible.

The thesis develops the problem in two methodological stages. First, it establishes a retained controller-centered rolling-horizon line, in which admission, protection, and risk-aware clipping decisions improve service performance under pressure. This line provides the main paper-facing experimental backbone and culminates in a disciplined controller progression under the `EXP00`, `EXP01`, and `Scenario1` settings. Second, the thesis shows through large-instance diagnostics that some remaining failures are better explained by sharply localized depot-wave bottlenecks than by broad admission pressure alone. This motivates a more integrated controller-routing-repair architecture implemented in a Julia-first fresh-solver framework.

The most important diagnostic result is that the West instance is dominated not by all-day capacity shortage, but by a morning picking wave concentrated around 05:30--06:15. To address this structure, the thesis introduces a dynamic shadow-price mechanism that converts bucket-level depot pressure into controller and routing signals. The resulting V8.2 variant achieves the best service--smoothing trade-off in the dynamic-price line: on West it reaches a service rate of approximately 99.41\% while reducing both total and morning picking overload relative to the earlier OR-V7 baseline. Additional validation shows that this mechanism also improves the medium-pressure Herlev instance, but is not universally dominant on the looser East instance, indicating clear regime dependence.

Overall, the thesis argues that rolling-horizon warehouse-distribution planning should be treated as an integrated control problem rather than as a route-only optimization task. A combination of controller-led admission logic, route-level feasibility management, and depot-aware smoothing provides a more realistic and more interpretable planning framework for complex multi-day delivery systems.
